\section{The finite-variance dictionary with Student hidden factors}
\label{sec:student-hidden-compression}
The normalized Gaussian readout need not be paired with Gaussian hidden factors. It is also useful with independent Student hidden factors of any fixed degrees of freedom $\nu>2$, retaining their logarithmic comparison cost for large teacher coefficients. The hidden gates are still independent $\operatorname{Gamma}(\kappa/N,1)$ variables. The final readout is still
\[
 z=G_0+\sum_j\sqrt{P_j}\,G_jh_j^H,
 \qquad P_j=S_j/\sum_iS_i,
\]
with standard Gaussian $G_0,G_j$ and a Dirichlet vector of total concentration $\kappa$. Only the hidden factors and hidden biases change to Student variables.

\begin{corollary}[Student-hidden normalized-readout compression]
\label{cor:student-hidden-compression}
For $q\ge1$, $0<\epsilon\le1/2$, and
\[
 L=H+q+1+\log\frac{d+1}{\epsilon},
\]
the native raw function is uniformly $\epsilon$-approximable, except on a prior event of probability at most $e^{-q}$, by a deterministic dictionary independent of $N$. Its networks have hidden width at most $CL$, coefficient bit length at most $CHL$, and
\[
 \log|\mathcal D|\le C[dHL^2+H^2L^3].
\]
Constants depend on the fixed activation, $\kappa$, and the Student degrees of freedom. If the hidden activation is bounded, the native raw supremum remains uniformly sub-Gaussian, independently of hidden depth and input dimension. The improved Gaussian-hidden $q^2\log q$ dictionary bound is not asserted for this model: the displayed budget is cubic in $q$ at fixed other parameters.
\end{corollary}
\begin{proof}
Finite first absolute moments of the Student factors give exactly the expected path and truncation recursions used in Theorem~\ref{thm:gaussian-ard-compression}. At the final layer, the normalized-readout proof uses only that the final proportions are independent of their own incoming features. Its deleted square-root mass is still $O(r)$ when coordinates with $P_j\le r^2$ are removed. Thus the same cutoff $r$ makes the expected uniform truncation error at most $e^{-q}\epsilon/8$.

The retained gate counts, the lower-tail control of the total final Gamma mass, and the internal gate cap are unchanged. Outside probability $O(e^{-q})$, each layer has at most $U=CL$ retained vertices and every internal gate is at most $CL$. Their retained factor arrays are selected using gates alone, so they retain their independent native laws. There are at most $M=C(dL+HL^2)$ hidden factors and biases. A Student tail union bound caps them by
\[
 B_G=1+(CM e^q)^{1/\nu},\qquad \log(B_G+1)\le CL.
\]
The Gaussian readout factors can be capped by $C\sqrt L$. Therefore the logarithm of every actual coefficient bound is $O(L)$. Deterministic layerwise Lipschitz propagation and dyadic rounding consequently require $CHL$ bits per coefficient. Enumerating the consecutive retained blocks gives
$CMHL\le C[dHL^2+H^2L^3]$ bits, as stated. All failure probabilities can be allocated as in the normalized-readout argument to sum to at most $e^{-q}$.

For a bounded hidden activation, the native supremum bound is
\[
 \|z\|_\infty\le|G_0|+M_\phi\sum_j\sqrt{P_j}|G_j|.
\]
Its Gamma--Dirichlet factorization depends only on final gates and readout factors. The previously proved $L^p$ bound $C\sqrt p$, and hence the sub-Gaussian tail, is unaffected by the hidden factor law.
\end{proof}
